% ==========================================
% SEZIONI DEL VERBALE
% ==========================================
\newpage
\vspace{0.8cm}

\section{Ordine del giorno}

\subsection{Gestione ruoli}
È stato deciso che fino al giorno 3 Novembre 2026, i ruoli sono definiti come segue:
\renewcommand{\labelitemii}{$\diamond$}
\begin{itemize}
    \item \textbf{redattori documento}:
    \begin{itemize}
        \item G. Gianmaria
        \item A. Tang
        \item F. Lissandrin
    \end{itemize}
    \item \textbf{verificatori}:
    \begin{itemize}
        \item C. Antoniazzo
        \item G. Gambarro
    \end{itemize}
    \item \textbf{amministratore}: T. Caon
    \item \textbf{responsabile}: D. Grosselle
\end{itemize}

\subsection{Template documenti}
È stato creato un template per la produzione dei documenti del gruppo per facilitare il lavoro degli amministratori da qui in avanti.

\subsection{Aggiornamento metodologia stesura documenti}
A seguito di un confronto interno, il gruppo ha rielaborato e formalizzato una nuova metodologia per la redazione, verifica e archiviazione della documentazione, al fine di garantire una maggiore tracciabilità.
\begin{enumerate}
    \item \textbf{Stesura della bozza iniziale (Redattori)}: I redattori designati per l'attività elaborano i contenuti preliminari su un documento collaborativo di Google Docs, garantendo la conclusione della prima bozza entro una tempistica utile a consentire le successive fasi di controllo;
    \item \textbf{Notifica di completamento}: Al termine della stesura, i redattori notificano tempestivamente la disponibilità della bozza sui canali di comunicazione interna del gruppo;
    \item \textbf{Fase di revisione (Verificatori)}: I verificatori procedono con l'ispezione del testo, annotando le richieste di modifica o correzione nella sezione commenti di Google Docs. Una volta terminato, oltre alla mail di sistema utilizzano il gruppo Whatsapp per informare i redattori;
    \item \textbf{Recepimento delle modifiche e aggiornamento del registro (Redattori)}: I redattori esaminano la sezione commenti, apportano le correzioni richieste e consolidano la versione del documento, curando l'aggiornamento del registro delle modifiche.
    \item \textbf{Trasposizione in LaTeX e archiviazione (Amministratore)}: Quando vi è riscontro positivo sulla versione consolidata, il documento viene integrato nel template LaTeX ufficiale del gruppo a cura dell'amministratore, procedendo alla compilazione del PDF e al caricamento nel repository di progetto e nello spazio di archiviazione condiviso, nell'apposita sottocartella.
\end{enumerate}

\subsection{Pianificazione riunioni}
Si stabilisce lo svolgimento di un meeting a cadenza settimanale ogni martedì pomeriggio fino all'aggiudicazione degli appalti, con eventuale passaggio a frequenza quindicinale per le fasi successive. Se necessario, verranno indetti meeting straordinari il venerdì per effettuare il punto della situazione in prossimità delle scadenze interne prefissate.
